% ===================================================================
%  அட்டவணை எண் 1 — பள்ளி, மேல்நிலைப் பள்ளி, வகுப்பறை.
%
%  Ceci n'est PAS une transcription : c'est une traduction, faite en
%  2026, en tamoul d'aujourd'hui. D'ou trois differences avec
%  texto/io et texto/fr, et trois seulement :
%
%   * pas de \nl ni de \cc — la traduction ne suit aucune ligne
%     imprimee, et les lignes de ce fichier ne sont que du confort ;
%   * pas de \begin{VUpage} — elle n'a ni feuillet ni folio ;
%   * le gras se pose sur le TERME seul, ponctuation dehors.
%
%  Les cles « %%K » sont celles de texto/io, macro pour macro pour
%  l'apparat, et les renvois \textsuperscript{(n)} visent les MEMES
%  objets de la planche, dans le MEME ordre.
%
%  DEUXIEME COLONNE DRAVIDIENNE, ET LE TELOUGOU N'EST PAS SA
%  VOISINE DE GRAPHIE. Les deux langues sont de la meme famille et
%  s'ecrivent dans deux alphabets qui n'ont pas une lettre en commun.
%  C'est l'inverse exact du couple hindi / marathi, qui partage
%  l'ecriture et non le lexique. La colonne tamoule ne risque donc
%  pas de derailler vers le telougou : elle risque autre chose.
%
%  SON ADVERSAIRE EST UN REGISTRE, ET C'EST LE PREMIER DU RELEVE. Le
%  tamoul est diglossique pour de bon : la langue parlee et la langue
%  ecrite different jusque dans la conjugaison — இருக்கு contre
%  இருக்கிறது, இல்ல contre இல்லை, பண்ணு contre செய் — et personne
%  n'imprime la premiere. Ecrire ce livret en tamoul parle serait
%  aussi faux que l'ecrire en argot : non pas trop moderne, mais d'un
%  autre usage. kolonoj.py releve donc six marques du parle, chacune
%  bornee pour ne pas mordre sur un mot honnete (இல்ல est aussi le
%  debut de இல்லை et de இல்லாத ; il faut exiger la fin de mot).
%
%  ET IL Y A UNE FAUTE SYMETRIQUE, QU'ON NE COMMETTRA PAS NON PLUS.
%  Le tamoul a un mouvement puriste vivant, et il serait facile de
%  remplacer chaque emprunt par une forge en தனித் தமிழ். Ce serait
%  fausser la colonne dans l'autre sens : les manuels de 2026
%  impriment நோட்டுப் புத்தகம் pour le cahier, சிலேட்டு pour
%  l'ardoise, பென்சில் pour le crayon. Ces mots-la sont ecrits, donc
%  ils s'ecrivent. On modernise la langue, jamais les choses — et on
%  ne la purifie pas non plus.
%
%  LES LETTRES GRANTHA NE SONT PAS UN DEFAUT, et c'est pourquoi le
%  controle ne les vise pas. ஜ, ஷ, ஸ, ஹ servent aux noms etrangers,
%  et ce tableau en est plein : ஸ்டெஃபானுஸ், ஜார்ஜியுஸ்,
%  ஃபிரான்சிஸ்குஸ். Un controle qui les refuserait interdirait a la
%  colonne d'ecrire les prenoms que la note (*) lui impose.
%
%  LA REGLE DU TELOUGOU VAUT ICI SANS RETOUCHE. Le tamoul est a tete
%  finale, son verbe finit la phrase, et il n'a pas de relative
%  postposee. parigi.py sert donc exactement comme pour le telougou
%  et le coreen ; c'est la quatrieme colonne a tete finale et il n'a
%  toujours pas fallu y toucher.
%
%  L'APPARAT PREND SES PROPRES MOTS : அட்டவணை pour un tableau — le
%  telougou voisin dit పట్టిక, et deux langues dravidiennes ne se
%  pretent pas plus leur apparat que deux langues indo-aryennes —,
%  தொடர் pour une serie, காட்சி pour une scene. காட்சி est le mot du
%  theatre tamoul : quatrieme colonne d'affilee a preferer le sien,
%  apres le प्रवेश marathe, le రంగం telougou et le 마당 coreen.
%
%  LA LIGNE « DI » EST LE SEUL ENDROIT OU L'APPARAT GENE. Le telougou
%  avait యొక్క, un genitif qui tient debout tout seul ; le tamoul n'en
%  a pas — son genitif est un suffixe. On ecrit உடைய, qui est un
%  participe et fait l'office au prix d'une lecture en apposition.
%  C'est la premiere fois qu'une macro de l'apparat resiste a une
%  langue, et il vaut mieux le dire que le masquer.
%
%  LA CLASSE TAMOULE A SON VOCABULAIRE, ET IL EST TAMOUL DE FOND :
%
%    கரும்பலகை (1)   le tableau noir — « la planche noire »
%    சுண்ணக்கட்டி (3)  la craie
%    மைக்கூடு (25)    l'encrier — « la maison de l'encre »
%    உறிஞ்சுத் தாள் (41) le buvard — « le papier qui aspire »
%    அகராதி (30)     le dictionnaire
%    கவராயம் (26)    le compas
%    அளவுகோல் (50)   la regle
%    மடிக்கத்தி (56)   le canif — « le couteau qui se plie »
%    குறிப்பேடு (54)   le carnet
%    பள்ளிப் பை (57)  le cartable
%    கடற்பஞ்சு (43)   l'eponge
%    இடைவேளை (76)   la recreation
%    பட்டறை (69)     l'atelier
%    மார்புச் சிலை (59) le buste
%    சிப்பந்தி (77)    l'homme de service
%
%  « சிப்பந்தி » ET LE « शिपाई » MARATHE SONT LE MEME MOT. Tous deux
%  viennent du persan sipāhī, tous deux ont cesse de designer un
%  soldat pour designer l'homme de service d'un bureau ou d'une
%  ecole. Cinquieme ecole reelle d'affilee — apres le فرّاش egyptien,
%  le शिपाई marathe, le బంట్రోతు telougou et le 수위 coreen — et la
%  premiere fois que deux colonnes tombent, sans se consulter, sur le
%  meme mot venu d'ailleurs.
%
%  ET LE VASISTAS (78) NE PASSE PAS EN UN MOT, POUR LA PREMIERE FOIS.
%  Trente-huit colonnes l'avaient nomme sans periphrase — le 들창
%  coreen, le గవాక్షం telougou, le postigi de l'ido lui-meme. Le
%  tamoul a la CHOSE dans toutes ses maisons, mais sa langue ecrite la
%  nomme en deux mots, மேல் சாளரம், la fenetre haute, et sa langue
%  parlee emprunte « ventilator » a l'anglais — emprunt que cette
%  colonne n'ecrit pas, parce qu'il n'est pas imprime. Le releve
%  disait « trente-huit fois sur trente-huit » ; il faut donc dire
%  maintenant trente-huit fois sur trente-neuf, et c'est la
%  trente-neuvieme qui est instructive : une serie qui ne casse
%  jamais ne prouve rien tant qu'on ne l'a pas cassee.
%
%  LES QUATRE OPERATIONS ET LES SCIENCES SONT DU TAMOUL PUR, sans un
%  emprunt : கூட்டல், கழித்தல், பெருக்கல், வகுத்தல் ; இயற்பியல்,
%  வேதியியல், தாவரவியல், விலங்கியல், நிலவியல், வரலாறு. La ou la
%  colonne emprunte, c'est pour les objets de la classe, pas pour ce
%  qu'on y enseigne — et cela se lit comme une date : le vocabulaire
%  savant a ete refait en tamoul, le mobilier scolaire ne l'a pas ete.
%
%  LA GAMME RESTE OCCIDENTALE, comme au telougou et au coreen. Le
%  tamoul a la sienne — ஸ ரி க ம ப த நி — et l'ecrire ici serait
%  tentant. Mais le tableau nomme do-re-mi et pose l'equivalence
%  « = C. D. E. F. G. A. B. » : c'est la gamme d'Europe qui est
%  l'objet, et l'on ne transpose pas les choses.
%
%  TROIS RETOURNEMENTS SUR LES HUIT PAIRES DE parigi.py, et ce sont
%  LES MEMES TROIS QU'AU COREEN : « stulo (8) en sua katedro (6) »,
%  « nigra tabelo (1) kun la vishilo (2) », « furnelo (46) kun lua
%  tubo (45) ». Les cinq autres sont exemptees par les trois motifs
%  connus — enumeration jointe par « e », renvoi de gauche deja
%  modificateur, frontiere de phrase. Que deux langues sans parente
%  tombent sur la meme liste n'est pas une coincidence : les trois
%  motifs sont des proprietes de l'ido, pas des langues d'arrivee.
% ===================================================================

%%K t01-apar-0 apar
\VUsaut{2.0mm}
\VUcentre{12.6pt}{120}{விளக்கச் சிறுநூல்}
\VUsaut{3.2mm}
\VUcentre{8.4pt}{160}{உடைய}
\VUsaut{2.6mm}
\VUtitre{15.6pt}{0}{டெல்மா - அட்டவணைகளுக்கு உதவும்}
\VUsaut{3.4mm}
\VUornamento{9pt}
\VUsaut{5.0mm}
\VUcentre{13.2pt}{90}{முதல் தொடர்}
\VUsaut{3.0mm}
\VUfilet{16mm}
\VUsaut{5.6mm}

%%K t01-apar-1 apar
\VUtitre{15.0pt}{60}{அட்டவணை எண் 1}
\VUsaut{1.4mm}
\VUtitre{11.4pt}{0}{பள்ளி, மேல்நிலைப் பள்ளி, வகுப்பறை.}
\VUsaut{2.2mm}
\VUfilet{20mm}
\VUsaut{6.4mm}

%%K t01-tit-1 sub
\VUpk{10.2pt}{40}{பொதுவான விளக்கம்.}
\VUsaut{3.0mm}

%%K t01-01-1 p
1. --- இந்த அட்டவணை ஒரு \VUgras{மேல்நிலைப் பள்ளியின்} \VUgras{வகுப்பறையைக்}
காட்டுகிறது. இந்த வகுப்பறையில் எட்டு \VUgras{மேசைகளும்} \textsuperscript{(18)}
எட்டு \VUgras{வாங்குகளும்} \textsuperscript{(32)} இருக்கின்றன. ஒவ்வொரு வாங்கிலும்
மூன்று மாணவர்கள் அமர்கிறார்கள். \VUgras{ஆசிரியர்} \textsuperscript{(7)} ஒரு
\VUgras{நாற்காலியில்} \textsuperscript{(8)} அமர்ந்திருக்கிறார் ; அந்த நாற்காலி
அவருடைய \VUgras{மேடையின்} \textsuperscript{(6)} மேல் இருக்கிறது.

%%K t01-02-1 p
2. --- மேடையின் இடப் பக்கத்தில் \VUgras{கண்ணாடி} போட்ட \VUgras{அலமாரி}
\textsuperscript{(15)} இருக்கிறது. வலப் பக்கத்தில் \VUgras{கரும்பலகை}
\textsuperscript{(1)} இருக்கிறது ; அதனுடன் \VUgras{துடைப்பான்}
\textsuperscript{(2)} இருக்கிறது, \VUgras{சுண்ணக்கட்டியும்}
\textsuperscript{(3)} இருக்கிறது.

%%K t01-02-2 p
மேடைக்கு மேலே \VUgras{சுவர்ப் படங்களும்} \textsuperscript{(9, 11,}
\textsuperscript{12)} \VUgras{வரைபடமும்} \textsuperscript{(10)} தொங்குகின்றன.

%%K t01-03-1 p
3. --- மாணவர்கள் எல்லாரும் தத்தம் \VUgras{இடத்தில்} \textsuperscript{(17)}
அமர்ந்திருக்கவில்லை. இயோவான்னெஸ் \textsuperscript{(4)} கரும்பலகையின் முன்
நின்று, \VUgras{துடைப்பானைப்} பிடித்துக் கொண்டு \VUgras{வடிவியல் உருவங்களை}
அழிக்கிறான்.

%%K t01-03-2 p
பேத்ருஸ் மேடைக்கு அருகில் இருக்கிறான். அவன் \VUgras{எழுத்துப் பணிகளை}
\textsuperscript{(5)} சேகரித்து ஆசிரியரிடம் கொண்டு போகிறான்.

%%K t01-04-1 p
4. --- \VUgras{இந்த} ஆசிரியர் \VUgras{நடைமுறை மொழிகளுக்கு} உரியவர். அவர்
மாணவர்களுக்கு அயல் மொழிகளைப் பேசவும் படிக்கவும் எழுதவும் கற்பிக்கிறார்
\VUgras{(ஜெர்மன், ஆங்கிலம், ஸ்பானிஷ், இத்தாலியன், ரஷ்யன்} அல்லது
\VUgras{பிரெஞ்சு)}.

%%K t01-05-1 p
5. --- வகுப்பறை மிகவும் பரந்ததாகவும் மிகவும் ஒளி மிக்கதாகவும் இருக்கிறது.
உள்பக்கத்தில் இரண்டு \VUgras{மேல் சாளரங்கள்} \textsuperscript{(78)} உள்ள ஒரு
அகன்ற \VUgras{சாளரமும்}, \VUgras{கண்ணாடிக் கதவும்} இருக்கின்றன ; அவை காற்று
வரவும் ஒளி வரவும் உதவுகின்றன.

%%K t01-05-2 p
இந்த வகுப்பறை குளிராக இல்லை. நடுவில், அறையைக் கதகதப்பாக்க மிகவும் நல்ல ஒரு
\VUgras{கணப்படுப்பு} \textsuperscript{(46)} இருக்கிறது ; அதற்கு ஒரு
\VUgras{புகைபோக்கி} \textsuperscript{(45)} உண்டு.

%%K t01-05-3 p
சுவரில் \VUgras{ஆடைக் கொக்கிகள்} \textsuperscript{(58)} பொருத்தப்பட்டிருக்கின்றன ;
அவற்றில் மாணவர்களின் தொப்பிகளும் மேலங்கிகளும் தொங்குகின்றன.

%%K t01-tit-2 sub
\VUsaut{5.2mm}
\VUpk{10.2pt}{40}{விளையாட்டு மைதானம்}
\VUsaut{3.6mm}

%%K t01-06-1 p
6. --- \VUgras{கடிகாரம்} \textsuperscript{(63)} ஒன்பது மணி ஐந்து நிமிடத்தைக்
காட்டுகிறது.

%%K t01-06-2 p
\VUgras{இடைவேளை} \textsuperscript{(76)} இப்போதுதான் முடிந்தது, பாடம் மீண்டும்
தொடங்குகிறது. விளையாடும் இடம் \VUgras{மைதானத்தில்} \textsuperscript{(75)}
இருக்கிறது. விளையாடும் மாணவர்கள் சிலரைப் பார்க்கிறோம்.
\VUgras{துணை ஆசிரியர்} \textsuperscript{(74)} அவர்களைக் கண்காணிக்கிறார்.

%%K t01-07-1 p
7. --- மைதானத்தில் வேறு பலரையும் பார்க்கிறோம், அதாவது \VUgras{ஆய்வாளர்}
\textsuperscript{(73)}, பள்ளித் தலைவர் \VUgras{(தலைமை ஆசிரியர்)}
\textsuperscript{(71)}, \VUgras{கண்காணிப்பாளர்} \textsuperscript{(72)}.

%%K t01-07-2 p
ஆய்வாளர் பள்ளிகளை ஆய்கிறார், தலைமை ஆசிரியர் மேல்நிலைப் பள்ளியை நடத்துகிறார்,
கண்காணிப்பாளர் படிப்பைக் கவனிக்கிறார்.

%%K t01-07-3 p
நான்காம் ஆள் \VUgras{சிப்பந்தி} \textsuperscript{(77)}. வராதவர்களைப் பதிவு
செய்வதற்காக அவர் அக்குளில் ஒரு பதிவேட்டை வைத்துக் கொண்டு போகிறார்.

%%K t01-08-1 p
8. --- மைதானத்தின் உள்பக்கத்தில் \VUgras{உடற்பயிற்சிக் கருவிகளும்}
\textsuperscript{(70)} \VUgras{கைவினைப் பணிகளுக்கான} \textsuperscript{(69)}
பட்டறையும் இருக்கின்றன.

%%K t01-08-2 p
அட்டவணையின் வலப் பக்கத்தில் \VUgras{இசை} \VUgras{வகுப்பறையும்}
\VUgras{ஓவிய} வகுப்பறையும் தெரியத் தொடங்குகின்றன.

%%K t01-noto-1 noto
\VUnotes{143.2mm}{%
(*) \textit{ஞானஸ்நானப் பெயர்களையும்} இலத்தீன் வடிவத்தின்படி எழுதுவது நல்லது.
எண்ணற்ற வேறுபாடுகளைத் தவிர்க்க அதுவே ஒரே வழி. மேலும் \textit{Giovanni, Jean,}
\textit{Johann, Jan, Hans, John, Ivan} என்பவற்றுள் எதைத் தேர்ந்தெடுப்பது?
ஞானஸ்நானப் பெயர்களுக்கென்று தனி அகராதி ஒன்றைச் செய்வோமா? இலத்தீனிலிருந்து
அவற்றின் எழுவாய் வடிவத்தை எடுத்து இப்படிச் சொல்வோம் : \VUgras{Ioannes, Ioanna;}
\VUgras{Iulius, Iulia; Iakobus; Andreas; Lukas.} (Gram. Kompleta).%
}

%%K t01-tit-3 sub
\VUpk{10.2pt}{40}{விரிவான விளக்கம்.}
\VUsaut{2.4mm}

%%K t01-tit-4 sub
\VUpk{10.2pt}{40}{நல்ல மாணவர்கள்.}
\VUsaut{3.0mm}

%%K t01-09-1 p
9. --- மேடையின் வலப் பக்கத்தில் உள்ள முதல் வாங்கின் மாணவர்கள்
\VUgras{ஹென்றிகுஸ்} (*), \VUgras{ஸ்டெஃபானுஸ்}, \VUgras{காரோலுஸ்}.

%%K t01-09-2 p
ஹென்றிகுஸ் மிகவும் கவனமாகவும் உழைப்பாகவும் இருக்கிறான், ஆசிரியர் சொல்வதைக்
கேட்கிறான்.

தன் மேசையின் மேல் அவனிடம் \VUgras{நோட்டுப் புத்தகம்} \textsuperscript{(28)},
\VUgras{புத்தகம்} \textsuperscript{(29)}, \VUgras{அகராதி} \textsuperscript{(30)},
\VUgras{எழுதுகோல் பெட்டி} \textsuperscript{(31)} இருக்கின்றன. அவனுடைய புத்தகத்தில்
\VUgras{பக்கங்கள்} பல, ஒவ்வொரு அதிகாரத்திலும் \VUgras{பத்திகள்} பல, ஒவ்வொரு
\VUgras{வரியிலும்} \VUgras{சொற்கள்} பல.

%%K t01-09-4 p
அவனுக்கு அடுத்திருக்கும் ஸ்டெஃபானுஸ் \textsuperscript{(24)} ஆர்வமுள்ளவன். அடுத்த
முறைக்கான பாடத்தைத் தன் நோட்டுப் புத்தகத்தில் குறித்து வைக்கிறான். அவனுக்கு
அழகான \VUgras{கையெழுத்து} உண்டு. அவனுடைய புத்தகம் மூடி இருக்கிறது. அதன்
\VUgras{அட்டையை} \textsuperscript{(27)} மட்டுமே பார்க்கிறோம். அவனுக்கு அருகில்
அவனுடைய \VUgras{கவராயம்} \textsuperscript{(26)} இருக்கிறது.

%%K t01-09-5 p
காரோலுஸ் \textsuperscript{(23)} தன் புத்தகத்தைக் கையால் பிடித்து, தன் பாடத்தை
மீண்டும் படிப்பதற்காக அதைத் திறக்கிறான். ஒவ்வொரு மாணவனையும் போலவே அவனுக்கும்
தன் முன்னால் \VUgras{மைக்கூடு} \textsuperscript{(25)} இருக்கிறது. அவன் சொல்
கேட்பவன்.

%%K t01-10-1 p
10. --- பக்கத்து மேசையில் \VUgras{லுடோவிகுஸ்}, \VUgras{அன்டோனியுஸ்},
\VUgras{பவுலுஸ்} இருக்கிறார்கள். லுடோவிகுஸ் தன் \VUgras{பாடத்தை}
\textsuperscript{(22)} ஒப்பிக்கிறான், ஆசிரியரின் \VUgras{வினாக்களுக்கு}
விடை சொல்கிறான். அன்டோனியுஸ் \textsuperscript{(21)} மிகவும் கவனமில்லாதவன்,
சில வேளைகளில் ஆசிரியர் அவனைத் தண்டிக்கவும் செய்கிறார். பவுலுஸ்
\textsuperscript{(20)} நல்ல \VUgras{தோழன்}, இனியவன், உதவி செய்பவன். அவன்
மிகவும் கவனமுள்ளவன்.

%%K t01-tit-5 sub
\VUsaut{4.2mm}
\VUpk{10.2pt}{40}{கெட்ட மாணவர்கள்.}
\VUsaut{3.2mm}

%%K t01-11-1 p
11. --- ஹென்றிகுஸுக்குப் பின்னால் \VUgras{ஜார்ஜியுஸ்} \textsuperscript{(38)}
இருக்கிறான். அவன் கெட்டவன் அல்ல, ஆனால் \VUgras{அளந்து பேசுபவன்},
கவனமில்லாதவன், குறும்புக்காரன். அவனுடைய \VUgras{லேசுத்தனமும்}
\VUgras{சொல் கேளாமையும்} பாடத்தின் போது \VUgras{கலவரத்தை} உண்டாக்குகின்றன,
அதனால் அடிக்கடி கடுமையான \VUgras{கண்டிப்புக்கு} ஆளாகிறான். அவனுடைய
\VUgras{கரடு ஏடு} \textsuperscript{(35)} அழுக்காக இருக்கிறது ; தன்
\VUgras{பேனாவால்} \textsuperscript{(34)} அதில் \VUgras{மையைச் சிந்துகிறான்},
அதனால் எப்போதும் தன் \VUgras{அழிப்பானை} \textsuperscript{(36)} உபயோகிக்கிறான் ;
அவனுடைய \VUgras{ஏட்டுப் பை} \textsuperscript{(33)} கிழிந்திருக்கிறது, ஏனெனில்
அவன் மிகவும் கவனக்குறைவானவன். நடைமுறை மொழிகளின் வகுப்பறைக்கு அவன் ஏன் தன்
\VUgras{பென்சில் பெட்டியைக்} \textsuperscript{(37)} கொண்டு வருகிறான்?

%%K t01-11-2 p
அவனுக்கு அடுத்திருக்கும் எட்முண்டுஸ் \textsuperscript{(39)} அவனை
விரும்புவதில்லை. அவன் உண்மையில் கொஞ்சம் சோம்பலானவன், ஆனால் பேச்சு
குறைந்தவன், அமைதியானவன். அவன் அளந்து பேசுவதில்லை ; ஆனால் கேட்பதற்குப் பதிலாகத்
தன் \VUgras{வாசிப்புப் புத்தகத்தின்} \VUgras{கதைகளைப்} படிக்கவே
விரும்புகிறான்.

%%K t01-11-3 p
இந்த வாங்கின் மூன்றாம் மாணவன் \VUgras{லுடோவிகுஸ்} \textsuperscript{(40)}. அவன்
உழைப்பாளி. வெள்ளைத் \VUgras{தாளில்} எழுதுகிறான். தன் முன்னால் தன்
\VUgras{உறிஞ்சுத் தாளை} \textsuperscript{(41)} வைத்திருக்கிறான்.

%%K t01-12-1 p
12. --- ஆசிரியரின் இடப் பக்கத்தில் உள்ள இரண்டாம் வாங்கின் மாணவர்கள் மிகவும்
சிறியவர்கள். முதலாமவன், சிறிய \VUgras{எமிலியுஸ்}, இன்னும் நன்றாக எழுதத்
தெரியாதவன் ; அவன் தன் \VUgras{சிலேட்டையும்} \textsuperscript{(42)}, தன்
\VUgras{சிலேட்டுக் குச்சியையும்}, தன் \VUgras{கடற்பஞ்சையும்}
\textsuperscript{(43)} கொண்டு வருகிறான்.

%%K t01-12-2 p
அவனுக்கு அருகில் \VUgras{அகுஸ்துஸ்}, \VUgras{பெர்னார்துஸ்}
\textsuperscript{(44)} இருக்கிறார்கள். இவன் நன்றாகப் படிக்கிறான், ஆனால்
அவனுடைய \VUgras{உடையணிதல்} மிகவும் கவனக்குறைவானது.

%%K t01-13-1 p
13. --- அவர்களுக்குப் பின்னால் மூன்று \VUgras{சோம்பேறிகள்} இருக்கிறார்கள்.
\VUgras{அல்பெர்துஸ்} தன் பாடங்களைக் கற்பதில்லை. \VUgras{யாகோபுஸ்}
\textsuperscript{(47)} கேட்பதற்குப் பதிலாகத் தூங்குகிறான். அவனுடைய
\VUgras{சோம்பல்} அவனுக்கு \VUgras{கடிந்துரைகளையும்} \VUgras{தண்டனைகளையும்}
கொண்டு வருகிறது. கடைசியாக \VUgras{கிறிஸ்தியானுஸ்} \textsuperscript{(48)}
மிகவும் நிதானமற்றவன். அவன் கெட்ட முறையில் \VUgras{நடந்து கொள்கிறான்},
திருந்துவதற்குக் கொஞ்சமும் முயல்வதில்லை.

%%K t01-14-1 p
14. --- அவனுக்கு மாறாக, அவனுடைய அடுத்தவர்கள் நல்ல நடத்தை உள்ளவர்கள்.
\VUgras{எர்னெஸ்துஸ்} \textsuperscript{(51)} ஒழுங்கையும் நேர்த்தியையும்
விரும்புகிறான், எப்போதும் தன் \VUgras{அளவுகோலையும்} \textsuperscript{(50)} தன்
\VUgras{பென்சிலையும்} \textsuperscript{(49)} உபயோகிக்கிறான். அவன்
\VUgras{பகல் மாணவன்}.

%%K t01-14-2 p
\VUgras{ஃபிரான்சிஸ்குஸ்} \textsuperscript{(52)} \VUgras{விடுதி மாணவன்} ; சீருடை
அணிந்து மேல்நிலைப் பள்ளியிலேயே சாப்பிட்டுத் தூங்குகிறான். அவன் நல்ல
\VUgras{தோழன்}.

%%K t01-14-3 p
மவுரிசியுஸ் தன் \VUgras{மடிக்கத்தியால்} \textsuperscript{(56)} தன்
\VUgras{பென்சிலைச்} சீவுகிறான். அவன் மிகவும் கவனமானவன்.

%%K t01-14-4 p
அவன் தன் புத்தகங்கள் எல்லாவற்றையும் கொண்டு வருகிறான் — தன்
\VUgras{இலக்கண நூலையும்} \textsuperscript{(55)}, தன் \VUgras{குறிப்பேட்டையும்}
\textsuperscript{(54)}, \VUgras{எழுதுத் தாங்கியையும்} \textsuperscript{(53)}
கூட. அவனுடைய \VUgras{பள்ளிப் பை} \textsuperscript{(57)} அவனுக்குப் பின்னால்
தரையில் இருக்கிறது.

%%K t01-14-5 p
வகுப்பறையின் உள்பக்கத்தில் உள்ள இரண்டு மேசைகளையும்
\VUgras{சிறந்த மாணவர்கள்} \textsuperscript{(19)} உபயோகிக்கிறார்கள்.

%%K t01-tit-6 sub
\VUsaut{4.6mm}
\VUpk{10.2pt}{40}{ஓவியமும் இசையும்.}
\VUsaut{3.4mm}

%%K t01-15-1 p
15. --- \VUgras{ஓவிய அறையில்} சுவரில் மாட்டப்பட்ட அழகான \VUgras{மாதிரிகளும்},
\VUgras{விளக்குக் கூடு} போட்ட ஒரு \VUgras{விளக்கும்} \textsuperscript{(62)}
கூரையிலிருந்து தொங்குகின்றன. \VUgras{ஆசிரியர்} \textsuperscript{(60)} மிகவும்
பொறுமையானவர். அவர் மாணவர்களின் \VUgras{ஓவியங்களைத்} \textsuperscript{(61)}
திருத்துகிறார், இயற்கையைப் பார்த்து \VUgras{மார்புச் சிலை}
\textsuperscript{(59)} வரையக் கற்றுத் தருகிறார்.

%%K t01-16-1 p
16. --- \VUgras{இசை அறை} மிகவும் சுவாரசியமாகத் தெரிகிறது. அங்கே
\VUgras{இசைக் குறிப்பைப்} படிக்கவும், \VUgras{ஐங்கோட்டின்} மேல் எழுதப்பட்ட
\VUgras{சுரக் குறிகளைப்} \textsuperscript{(67)} பாடவும் (த, ரி, க, ம, ப, த,
நி\,=\,C. D. E. F. G. A. B.), \VUgras{சுரக் குறியீடுகளை} அறியவும், அழகான
\VUgras{இன்னிசைகளைப்} பாடவும் கற்றுக் கொள்கிறார்கள். இசைத் தாள்களை
\VUgras{இசைத்தாள் தாங்கியின்} \textsuperscript{(65)} மேல் வைக்கிறார்கள்.
மாணவர்களுள் ஒருவன் \VUgras{பியானோ வாசிக்கிறான்} \textsuperscript{(64)},
\VUgras{இசை ஆசிரியர்} \textsuperscript{(66)} \VUgras{தாளம் போடுகிறார்}
\textsuperscript{(68)}.

%%K t01-17-1 p
17. --- \VUgras{கைவினைப் பணிகளுக்கான} \textsuperscript{(69)}
\VUgras{பட்டறையின்} ஒரு மூலை தெரியத் தொடங்குகிறது ; அங்கே சிறுவர்கள் மரத்தைச்
சீவவும் இரும்பை அராவவும் கற்கலாம். எல்லா மேல்நிலைப் பள்ளிகளிலும் இது
இருப்பதில்லை.

%%K t01-tit-7 sub
\VUsaut{5.0mm}
\VUpk{10.2pt}{40}{அறிவியல் வகுப்பறை.}
\VUsaut{3.6mm}

%%K t01-18-1 p
18. --- கணக்கு ஆசிரியர் மாணவர்களுக்கு \VUgras{வடிவியல்},
\VUgras{எண்கணிதம்}, \VUgras{கணக்கு}, \VUgras{மெட்ரிக் முறை} ஆகியவற்றைக்
கற்பிக்கிறார். கரும்பலகையின் மேல் முதன்மையான வடிவியல் உருவங்களைப் பார்க்கிறோம் :
\VUgras{வட்டம்} \textsuperscript{(a)}, \VUgras{சதுரம்} \textsuperscript{(b)},
\VUgras{முக்கோணம்} \textsuperscript{(c)}, \VUgras{கோடு} \textsuperscript{(d)}.

%%K t01-18-2 p
ஒரு \VUgras{கணக்கையும்} பார்க்கிறோம். அது \VUgras{கூட்டலும்} அல்ல,
\VUgras{கழித்தலும்} அல்ல, \VUgras{பெருக்கலும்} அல்ல ; அது
\VUgras{வகுத்தல்}தான் \textsuperscript{(e)}.

%%K t01-19-1 p
19. --- மெட்ரிக் முறையின் அட்டவணையின் மேல் இவற்றைப் பார்க்கிறோம் :
\VUgras{மீட்டர்} \textsuperscript{(a)}, \VUgras{தராசு} \textsuperscript{(b)}
அதனுடைய \VUgras{எடைக் கற்கள்}, \VUgras{லிட்டர்} \textsuperscript{(e)},
\VUgras{டெக்காலிட்டர்} \textsuperscript{(f)}, \VUgras{டெசிலிட்டர்},
\VUgras{கன மீட்டர்} \textsuperscript{(d)}.

%%K t01-19-2 p
மாணவர்கள் \VUgras{இயற்கை அறிவியலையும்} \textsuperscript{(11)}, விலங்கியலையும்
\textsuperscript{(a)}, தாவரவியலையும் \textsuperscript{(b)}, நிலவியலையும்
\textsuperscript{(c)}, \VUgras{வரலாற்றையும்} \textsuperscript{(12)}
படிக்கிறார்கள்.

%%K t01-19-3 p
அலமாரிக்குள் \VUgras{இயற்பியல்} \textsuperscript{(15)},
\VUgras{வேதியியல்} \textsuperscript{(16)} கருவிகள் இருக்கின்றன.

%%K t01-19-4 p
அலமாரியின் மேல் \VUgras{கோளம்} \textsuperscript{(14)}, \VUgras{கூம்பு},
\VUgras{நிலவுருண்டை} \textsuperscript{(13)} இருக்கின்றன.

%%K t01-19-5 p
அட்டவணைகளுக்கு மேலே உலகின் ஐந்து பகுதிகளைக் காட்டும் \VUgras{வரைபடம்}
தொங்குகிறது : \VUgras{ஐரோப்பா} \textsuperscript{(g)}, \VUgras{ஆசியா}
\textsuperscript{(e)}, \VUgras{ஆப்பிரிக்கா} \textsuperscript{(f)},
\VUgras{அமெரிக்கா} \textsuperscript{(ab)}, \VUgras{ஓசியானியா}
\textsuperscript{(h)}, மேலும் இரண்டு பெரிய கடல்களான \VUgras{பசிபிக்}
\textsuperscript{(c)}, \VUgras{அட்லாண்டிக்} \textsuperscript{(d)}.

\VUsaut{6.0mm}
\VUfilet{22mm}
